% Created 2026-05-11 Mon 02:18
% Intended LaTeX compiler: lualatex
\documentclass[11pt]{article}

\usepackage{amsmath}
\usepackage{fontspec}
\usepackage{graphicx}
\usepackage{longtable}
\usepackage{wrapfig}
\usepackage{rotating}
\usepackage[normalem]{ulem}
\usepackage{capt-of}
\usepackage{hyperref}
\usepackage{minted}
\usepackage{fontspec}
\setmainfont{UT Sans}[
BoldFont={UT Sans Bold},
UprightFont={UT Sans Regular},
FontFace={m}{n}{UT Sans Medium}
]
% Fallback fake italic using a slant
\usepackage{etoolbox}
\newcommand{\fakeitalic}[1]{{\addfontfeatures{FakeSlant=0.18}#1}}
\AtBeginDocument{\renewcommand{\emph}[1]{\fakeitalic{#1}}}
\usepackage{minted}
\usepackage[a4paper,margin=3cm]{geometry}
\usepackage{lastpage}
\usepackage{fancyhdr}
\usepackage{lastpage}
\pagestyle{fancy}                % set fancy style
\fancyhf{}                        % clear all headers and footers
\cfoot{\thepage/\pageref*{LastPage}}  % center footer
\renewcommand{\headrulewidth}{0pt}    % remove header line
\renewcommand{\footrulewidth}{0pt}    % remove footer line
% also override plain page style for first page
\fancypagestyle{plain}{%
\fancyhf{}%
\cfoot{\thepage/\pageref*{LastPage}}%
\renewcommand{\headrulewidth}{0pt}%
\renewcommand{\footrulewidth}{0pt}%
}
\usepackage{url}
\author{Elias-Robert BAJCSI}
\date{}
\title{Laborator 3\\\medskip
\large PLFP}
\hypersetup{
 pdfauthor={Elias-Robert BAJCSI},
 pdftitle={Laborator 3},
 pdfkeywords={},
 pdfsubject={},
 pdfcreator={},
 pdflang={English}}
\begin{document}

\maketitle
\section{\texttt{./Cargo.toml}}
\label{sec:org43d9b33}
\begin{minted}[breaklines=true,breakanywhere=true,breaksymbol=,fontsize=\small,linenos=false]{toml}
[package]
name = "lab3"
version = "0.1.0"
edition = "2021"

[[bin]]
name = "ex1"
path = "src/bin/ex1.rs"

[[bin]]
name = "ex2"
path = "src/bin/ex2.rs"

[[bin]]
name = "ex3"
path = "src/bin/ex3.rs"

[[bin]]
name = "ex4"
path = "src/bin/ex4.rs"
\end{minted}
\section{\texttt{./src/bin/ex1.rs}}
\label{sec:org093fbb5}
\begin{minted}[breaklines=true,breakanywhere=true,breaksymbol=,fontsize=\small,linenos=false]{rs}
fn cel_mai_mare(x: i32, y: i32) -> i32 {
    if x > y { x } else { y }
}

fn cel_mai_mic(x: i32, y: i32) -> i32 {
    if x < y { x } else { y }
}

fn alege_comparator(crescator: bool) -> fn(i32, i32) -> i32 {
    if crescator {
        cel_mai_mare
    } else {
        cel_mai_mic
    }
}

fn main() {
    let comparatoare: Vec<fn(i32, i32) -> i32> = vec![cel_mai_mare, cel_mai_mic];
    let perechi = [(3, 7), (10, 2), (5, 5)];

    for cmp in &comparatoare {
        for (x, y) in perechi {
            println!("({}, {}) -> {}", x, y, cmp(x, y));
        }
    }

    let cmp1 = alege_comparator(true);
    let cmp2 = alege_comparator(false);

    println!("true  -> {}", cmp1(10, 2));
    println!("false -> {}", cmp2(10, 2));
}
\end{minted}
\section{\texttt{./src/bin/ex2.rs}}
\label{sec:org5648cb6}
\begin{minted}[breaklines=true,breakanywhere=true,breaksymbol=,fontsize=\small,linenos=false]{rs}
fn aplica_tuturor(lista: &[i32], op: impl Fn(i32) -> i32) -> Vec<i32> {
    lista.iter().map(|x| op(*x)).collect()
}

fn main() {
    let n: i32 = 5;
    let inmulteste_cu_n = |x: i32| x * n;

    for x in 1..=10 {
        println!("{}", inmulteste_cu_n(x));
    }

    let lista: Vec<i32> = (1..=10).collect();
    let rezultat = aplica_tuturor(&lista, inmulteste_cu_n);

    println!("{:?}", rezultat);
}
\end{minted}
\section{\texttt{./src/bin/ex3.rs}}
\label{sec:orgc124e28}
\begin{minted}[breaklines=true,breakanywhere=true,breaksymbol=,fontsize=\small,linenos=false]{rs}
struct Data {
    zi: u32,
    luna: u32,
    an: u32,
}

fn fmt_literar(d: &Data) -> String {
    let luni = ["Jan", "Feb", "Mar", "Apr", "May", "Jun", "Jul", "Aug", "Sep", "Oct", "Nov", "Dec"];
    format!("{} {} {}", d.zi, luni[(d.luna - 1) as usize], d.an)
}

fn fmt_iso(d: &Data) -> String {
    format!("{:04}-{:02}-{:02}", d.an, d.luna, d.zi)
}

fn fmt_slash_dmy(d: &Data) -> String {
    format!("{:02}/{:02}/{:04}", d.zi, d.luna, d.an)
}

fn fmt_slash_mdy(d: &Data) -> String {
    format!("{:02}/{:02}/{:04}", d.luna, d.zi, d.an)
}

fn fmt_punct(d: &Data) -> String {
    format!("{:04}.{:02}.{:02}", d.an, d.luna, d.zi)
}

fn fmt_scurt(d: &Data) -> String {
    format!("{}/{}/{}", d.zi, d.luna, d.an % 100)
}

fn fmt_literar_zero(d: &Data) -> String {
    let luni = ["Jan", "Feb", "Mar", "Apr", "May", "Jun", "Jul", "Aug", "Sep", "Oct", "Nov", "Dec"];
    format!("{:02} {} {}", d.zi, luni[(d.luna - 1) as usize], d.an)
}

fn adauga_unu(x: i32) -> i32 { x + 1 }
fn dublu(x: i32) -> i32 { x * 2 }
fn scade_trei(x: i32) -> i32 { x - 3 }
fn patrat(x: i32) -> i32 { x * x }

fn aplica_pipeline(val: i32, pasi: &[fn(i32) -> i32]) -> i32 {
    let mut rezultat = val;
    for (i, f) in pasi.iter().enumerate() {
        rezultat = f(rezultat);
        println!("pas {} -> {}", i + 1, rezultat);
    }
    rezultat
}

fn main() {
    let data = Data { zi: 5, luna: 4, an: 1981 };

    let formate: Vec<fn(&Data) -> String> = vec![
        fmt_literar,
        fmt_iso,
        fmt_slash_dmy,
        fmt_slash_mdy,
        fmt_punct,
        fmt_scurt,
        fmt_literar_zero,
    ];

    for fmt in &formate {
        println!("{}", fmt(&data));
    }

    let pasi: Vec<fn(i32) -> i32> = vec![adauga_unu, dublu, scade_trei, patrat];
    let rezultat = aplica_pipeline(5, &pasi);
    println!("rezultat final: {}", rezultat);
}
\end{minted}
\section{\texttt{./src/bin/ex4.rs}}
\label{sec:orgfcc3500}
\begin{minted}[breaklines=true,breakanywhere=true,breaksymbol=,fontsize=\small,linenos=false]{rs}
#[derive(Debug, Clone)]
struct Fel {
    id: u32,
    descriere: String,
    calorii: u32,
    pret: f64,
}

impl Fel {
    fn nou(id: u32, descriere: &str, calorii: u32, pret: f64) -> Fel {
        Fel {
            id,
            descriere: descriere.to_string(),
            calorii,
            pret,
        }
    }
}

fn afis_scurt(f: &Fel) -> String {
    format!("[{}] {}", f.id, f.descriere)
}

fn afis_complet(f: &Fel) -> String {
    format!("[{}] {} | {} kcal | {:.2} lei", f.id, f.descriere, f.calorii, f.pret)
}

fn tipareste_meniu(mese: &[Fel], fmt: fn(&Fel) -> String) {
    for f in mese {
        println!("{}", fmt(f));
    }
}

fn main() {
    let mese = vec![
        Fel::nou(1, "Supa de legume", 210, 17.0),
        Fel::nou(2, "Piept de pui cu orez", 540, 31.5),
        Fel::nou(3, "Clatite", 430, 19.0),
    ];

    let formatari: Vec<fn(&Fel) -> String> = vec![afis_scurt, afis_complet];

    for fmt in &formatari {
        for f in &mese {
            println!("{}", fmt(f));
        }
        println!();
    }

    tipareste_meniu(&mese, afis_scurt);
    println!();
    tipareste_meniu(&mese, afis_complet);
}
\end{minted}
\end{document}
